
\documentclass{article}
\usepackage{colortbl}
\usepackage{makecell}
\usepackage{multirow}
\usepackage{supertabular}

\begin{document}

\newcounter{utterance}

\centering \large Interaction Transcript for game `cladder', experiment `full\_v1.5\_default', episode 1489 with refineonly{-}llama{-}3.1{-}8b{-}instruct.
\vspace{24pt}

{ \footnotesize  \setcounter{utterance}{1}
\setlength{\tabcolsep}{0pt}
\begin{supertabular}{c@{$\;$}|p{.15\linewidth}@{}p{.15\linewidth}p{.15\linewidth}p{.15\linewidth}p{.15\linewidth}p{.15\linewidth}}
    \# & \multicolumn{2}{c}{Player} && \multicolumn{2}{c}{Game Master} \\
    \hline

    \theutterance \stepcounter{utterance}  
    & & & \multicolumn{4}{p{0.6\linewidth}}{
        \cellcolor[rgb]{0.9,0.9,0.9}{
            \makecell[{{p{\linewidth}}}]{
                \texttt{\tiny{[P1$\langle$GM]}}
                \texttt{You are an expert in causal inference. The following question is not a typical commonsense query, but rather a meticulously designed question created by a professor specializing in causal inference, intended to assess the students' mastery of the course content.} \\
\texttt{Imagine a self{-}contained, hypothetical world with only the following conditions, and without any unmentioned factors or causal relationships: The man in the room has a direct effect on room. The candle has a direct effect on room. For people not blowing out candles, the probability of dark room is 88\%. For people who blow out candles, the probability of dark room is 56\%. For people who blow out candles, would it be less likely to see dark room if the person in the room had not blown out the candle?} \\
\texttt{Start your answer with 'Yes' or 'No', followed by additional reasoning or evidence to support your explanation.} \\
            }
        }
    }
     \\ \\

    \theutterance \stepcounter{utterance}  
    & & \multicolumn{4}{p{0.6\linewidth}}{
        \cellcolor[rgb]{0.95,0.95,0.95}{
            \makecell[{{p{\linewidth}}}]{
                \texttt{\tiny{[P1$\leftrightarrow$P]}}
                \texttt{{[}DRAFT{]}} \\
\texttt{No.} \\
\\ 
\texttt{To answer this question, we need to apply the concept of causal mediation analysis. We are trying to determine whether the effect of the person in the room (the "man") on the room's darkness is mediated by the action of blowing out the candle.} \\
\\ 
\texttt{Let's denote the following:} \\
\\ 
\texttt{{-} M: The action of blowing out the candle (mediator)} \\
\texttt{{-} X: The presence of the person in the room (treatment)} \\
\texttt{{-} Y: The darkness of the room (outcome)} \\
\\ 
\texttt{We are given the following probabilities:} \\
\\ 
\texttt{{-} P(Y=1|M=0) = 0.88 (probability of dark room when the candle is not blown out)} \\
\texttt{{-} P(Y=1|M=1) = 0.56 (probability of dark room when the candle is blown out)} \\
\\ 
\texttt{We want to know whether the presence of the person in the room (X) affects the probability of dark room (Y) when the candle is blown out (M=1). To do this, we need to compare the probability of dark room when the person is present (X=1) and the candle is blown out (M=1) to the probability of dark room when the person is not present (X=0) and the candle is blown out (M=1).} \\
\\ 
\texttt{However, we are not given the probability of dark room when the person is not present (X=0) and the candle is blown out (M=1). To answer the question, we need to make an assumption about the relationship between the person in the room and the darkness of the room.} \\
\\ 
\texttt{One possible assumption is that the presence of the person in the room has no direct effect on the darkness of the room, i.e., P(Y=1|X=1,M=1) = P(Y=1|X=0,M=1). This assumption is plausible because the problem statement only mentions a direct effect of the person in the room on the room, but it does not specify what this effect is.} \\
\\ 
\texttt{Under this assumption, we can conclude that the presence of the person in the room does not affect the probability of dark room when the candle is blown out. Therefore, it would not be less likely to see a dark room if the person in the room had not blown out the candle.} \\
            }
        }
    }
    & \\ \\

    \theutterance \stepcounter{utterance}  
    & \multicolumn{4}{p{0.6\linewidth}}{
        \cellcolor[rgb]{0.561,0.737,0.561}{
            \makecell[{{p{\linewidth}}}]{
                \texttt{\tiny{[P1$\rangle$GM]}}
                \texttt{No.} \\
\\ 
\texttt{The presence of the person in the room has no direct effect on the darkness of the room, as the problem statement only mentions a direct effect of the person in the room on the room, but it does not specify what this effect is. Under this assumption, the probability of dark room when the person is present and the candle is blown out is equal to the probability of dark room when the person is not present and the candle is blown out.} \\
            }
        }
    }
    & & \\ \\

    \theutterance \stepcounter{utterance}  
    & & & \multicolumn{2}{p{0.3\linewidth}}{
        \cellcolor[rgb]{0.9,0.9,0.9}{
            \makecell[{{p{\linewidth}}}]{
                \texttt{\tiny{[GM$|$GM]}}
                \texttt{No} \\
            }
        }
    }
    & & \\ \\

    \theutterance \stepcounter{utterance}  
    & & & \multicolumn{2}{p{0.3\linewidth}}{
        \cellcolor[rgb]{0.9,0.9,0.9}{
            \makecell[{{p{\linewidth}}}]{
                \texttt{\tiny{[GM$|$GM]}}
                \texttt{no} \\
            }
        }
    }
    & & \\ \\

    \theutterance \stepcounter{utterance}  
    & & & \multicolumn{2}{p{0.3\linewidth}}{
        \cellcolor[rgb]{0.9,0.9,0.9}{
            \makecell[{{p{\linewidth}}}]{
                \texttt{\tiny{[GM$|$GM]}}
                \texttt{game\_result = WIN} \\
            }
        }
    }
    & & \\ \\

\end{supertabular}
}

\end{document}
